% ch3.tex


\chapter{প্রাচীন ব্যাবিলনীয় নক্ষত্রবিদ্যা}
\epigraph{``\textit{
খেলিছ এ বিশ্ব লয়ে \\
বিরাট শিশু আনমনে। \\
প্রলয় সৃষ্টি তব পুতুল খেলা \\
নিরজনে প্রভু নিরজনে।। \\
}''}{\textit{নজরুল ইসলাম}}

যে প্রাচীনতম কাল থেকে তথ্য আমাদের কাছে পৌঁছেছে, তখন ইউফ্রেটিস ও টাইগ্রিসের মধ্যবর্তী সমভূমি নিয়ে গঠিত শিনার দেশে দুটি স্বতন্ত্র জাতির বাস ছিল। উত্তরাংশ আক্কাদে ছিল একটি সেমিটিক জাতি, আর দক্ষিণাংশে সুমেরীয়রা। ভাষা, চেহারা ও রীতিনীতিতে এই দুই জাতি একেবারেই আলাদা ছিল। সুমেরীয়রা ছিল আদি অধিবাসী বা প্রথম অভিবাসী। অন্যদিকে সেমিটিরা পরে পশ্চিমের মরুভূমি থেকে ক্রমবর্ধমান সংখ্যায় এখানে আসে। সুমেরীয় ভাষার সঙ্গে সেমিটিক বা ইন্দো-ইউরোপীয় ভাষার কোনো সম্পর্ক দেখা যায় না, বরং তুরানীয় ভাষার সঙ্গে সাদৃশ্য আছে। যা-ই হোক, এই প্রাচীনতম সংস্কৃতির জাতির নৃতাত্ত্বিক পরিচয় অজানা।

খ্রিস্টপূর্ব ৩০০০ সালের পরের শতাব্দীগুলোতে দক্ষিণের সুমেরীয় নগরগুলোতে সংস্কৃতি সবচেয়ে উন্নত ছিল: সমুদ্রতীরের এরিদু, মরুভূমির প্রান্তে উর, এবং উরুক, লাগাশ, নিপ্পুর ও লারসা। সাধারণত এগুলোর একটির হাতে আধিপত্য থাকত। উত্তরের নগরগুলোর (আগাদে, সিপ্পার, বোরসিপ্পা, ব্যাবিলন) সেমিটিরা এই সংস্কৃতি গ্রহণ করে, এবং আরও অভিবাসনের পর ক্রমশ প্রধান জাতিতে পরিণত হয়। খ্রিস্টপূর্ব ২৫০০ সালের দিকে আগাদের সারগন এবং পরে তাঁর পুত্র নারাম-সিন যখন সমগ্র মেসোপটেমিয়া শাসন করতেন, তখন সামরিক অংশ ছিল সেমিটিদের নিয়ে গঠিত, আর লিপিকার ও বেসামরিক কর্মকর্তারা ছিলেন সুমেরীয়। পরের শতাব্দীগুলোতে কর্তৃত্ব আবার দক্ষিণে ফিরে যায়। সেখানে লাগাশ ও উরের শাসকেরা নিজেদের "সুমের ও আক্কাদের রাজা" বলতেন। খ্রিস্টপূর্ব ২০০০ সালের পর পশ্চিম থেকে আমোরীয় গোত্রগুলোর অভিবাসনের ফলে সেমিটিরা নিশ্চিত প্রাধান্য লাভ করে। তখন একটি রাজবংশের অধীনে, যার সবচেয়ে বিখ্যাত রাজা হাম্মুরাবি, ব্যাবিলন নগরী একটি বিশাল সাম্রাজ্যের রাজধানী এবং বাণিজ্য ও সংস্কৃতির কেন্দ্র হয়ে ওঠে।

সুমেরীয়রাই কীলকাকার লিপির উদ্ভাবক। এই লিপিতে একটি স্বরবর্ণ এবং এক বা দুটি ব্যঞ্জনবর্ণ নিয়ে গঠিত প্রতিটি ধ্বনির জন্য একটি বিশেষ চিহ্ন ছিল। নরম কাদার ফলকে একটি লেখনী, অর্থাৎ ধাতব কীলক, দিয়ে ছাপ দিয়ে চিহ্নগুলো তৈরি করা হতো। ছাপগুলো এক পাশে চওড়া, অন্য পাশে সরু হতো। তারপর ফলকটি আগুনে পুড়িয়ে শক্ত করা হতো। চিহ্ন দিয়ে অক্ষর প্রকাশের এই পদ্ধতি সুমেরীয়ের মতো একটি সংযোজনশীল (agglutinating) ভাষার জন্য পুরোপুরি উপযোগী ছিল, যেখানে আলাদা আলাদা শব্দ ও ধাতু কেবল পরপর জুড়ে দেওয়া হতো। সেমিটিরা এই কীলকাকার লিপি গ্রহণ করে, যদিও ধাতুর রূপান্তরভিত্তিক তাদের ভাষার জন্য এটি বেশ অসুবিধাজনক ছিল। ব্যাবিলন যখন সমগ্র নিকট প্রাচ্যের সাংস্কৃতিক কেন্দ্র হয়ে উঠল, তখন এর ভাষা ও কীলকাকার লিপি মিশর ও এশিয়া মাইনর পর্যন্ত আন্তর্জাতিক যোগাযোগে ব্যবহৃত হতো।

কীলকাকার লিপি এবং তাতে লেখা ভাষাগুলোর পাঠোদ্ধার ছিল উনিশ শতকের অন্যতম বড় কৃতিত্ব। এর ফলে ইতিহাস ও সংস্কৃতির একটি সম্পূর্ণ হারিয়ে যাওয়া যুগ আলোয় এল। আগের জ্ঞান বেশিরভাগই এসেছিল গ্রিক উৎস থেকে, যা ছিল খণ্ডিত ও উপাখ্যাননির্ভর, এবং তা খ্রিস্টপূর্ব ৭০০ সালের আগে যেত না। ১৮৪৬ সালে ফরাসি কনসাল বোত্তার প্রথম খননকাজে অনুপ্রাণিত হয়ে হেনরি লেয়ার্ড প্রাচীন নিনেভের স্থানে খনন শুরু করেন। তিনি অপূর্ব সব ভাস্কর্য খুঁড়ে বের করেন: শিকারের দৃশ্যের উদ্‌গত খোদাই, ডানাওয়ালা ষাঁড় ও ড্রাগন। চারপাশে ছড়িয়ে থাকা কীলকাকার লিপিযুক্ত অসংখ্য ভাঙা টুকরো দেখে তিনি বিস্মিত হন। সেগুলোর মূল্য আন্দাজ করে তিনি শিল্পকর্মগুলোর সঙ্গে বাক্সভর্তি এই টুকরোও ব্রিটিশ মিউজিয়ামে পাঠান। বহু বছর পর জর্জ স্মিথ যখন মহাপ্লাবনের কাহিনিসংবলিত কিছু লেখার পাঠোদ্ধার করেন, তখন এগুলোর গুরুত্ব বোঝা যায়। তারপর যত বেশি সম্ভব টুকরো সংগ্রহের জন্য বিশেষ অভিযান পাঠানো হয়। দেখা গেল, এগুলো রাজা আশুরবানিপালের "গ্রন্থাগারের" অবশেষ। এতে নতুন শিলালিপির পাশাপাশি আগের শতাব্দীগুলোর লেখার অনুলিপি, অভিধান ও অভিধান-সংক্রান্ত উপকরণ ছিল। এগুলো প্রাচীন মেসোপটেমীয় সংস্কৃতির চর্চায় (যাকে তখন থেকে "অ্যাসিরিওলজি" বলা হয়) জোরালো প্রেরণা জোগায়।

শুরুতে দুটি সম্পূর্ণ ভিন্ন ভাষার মিশ্রণ অ্যাসিরিওলজিস্টদের বিভ্রান্ত করেছিল। সাধারণত প্রতিটি চিহ্ন ধ্বনি হিসেবে ব্যবহৃত হতো, অর্থাৎ সেমিটিক ভাষার জন্য সুমেরীয় থেকে একটি অক্ষর হিসেবে ধার নেওয়া হতো। কিন্তু প্রায়ই চিহ্নটি সুমেরীয় ভাষায় তার অর্থ অনুযায়ী তথাকথিত "ভাবচিত্র" (ideogram) হিসেবে, অর্থাৎ কোনো বস্তু বা ধারণার প্রতীক হিসেবে ব্যবহৃত হতো। এভাবে দুটি ভাষারই রহস্য ধীরে ধীরে একসঙ্গে উন্মোচিত হলো। অন্যান্য প্রাচীন নগরের ধ্বংসাবশেষ থেকে খুঁড়ে পাওয়া লেখা এই সাহিত্যিক খণ্ডগুলোকে পূর্ণতা দিল, আর এগুলো প্রাচীনকালের ইতিহাস ও সংস্কৃতিকে ক্রমশ আরও স্পষ্ট করে তুলল।

এর মাধ্যমে জ্যোতির্বিজ্ঞানের একটি অজানা প্রাচীন ইতিহাসও আলোয় এল। এই পুনরাবিষ্কৃত প্রাচীন জগতে জ্যোতিষ্কগুলো অন্য যেকোনো দেশ বা যুগের চেয়ে বেশি গুরুত্ব পেয়েছিল। ছোট ছোট বিক্ষিপ্ত খণ্ড থেকে এই প্রাচীন জ্ঞান পুনর্গঠন করা অবশ্যই কঠিন ছিল। অস্পষ্ট অনুমান, ঝুঁকিপূর্ণ ধারণা, ভুল ব্যাখ্যা আর অসমর্থনযোগ্য তত্ত্বের মধ্য দিয়েই কেবল ধীরে ধীরে সত্যের কাছাকাছি পৌঁছানো গেছে। বহু বছর ধরে হুগো উইংক্লারের তত্ত্ব, তথাকথিত "প্যান-ব্যাবিলনবাদ", খুব খ্যাতি পেয়েছিল। এই তত্ত্ব ঘোষণা করত যে প্রাচীনতম কালে, খ্রিস্টপূর্ব ৩০০০ থেকে ২০০০ সালের মধ্যে, সেখানে ইতিমধ্যেই একটি অত্যন্ত উন্নত জ্যোতির্বিজ্ঞান ছিল। এর ভিত্তি ছিল আকাশের পর্যায়কালগুলোর পূর্ণাঙ্গ জ্ঞান, এবং বিষুবের অয়নচলনের (precession of the equinoxes) ফলে তারামণ্ডলের চেহারা কীভাবে সরে যায় তার জ্ঞান। তত্ত্বটি দাবি করত যে এই আদিম বিশ্বধারণা পার্থিব ও আকাশের ঘটনাবলির মধ্যে ঘনিষ্ঠ সম্পর্ক ঘোষণা করত বলে এর চরিত্র ছিল "নাক্ষত্রিক পুরাণের"। এটিই পরবর্তী সব প্রাচ্য ও গ্রিক চিন্তাধারার উৎস, এবং আধুনিক কাল পর্যন্ত ইউরোপের দূরবর্তী অঞ্চলেও কিংবদন্তি, লোকজ্ঞান ও রীতিনীতি নির্ধারণ করেছে। পরে লেখাগুলো যত্নের সঙ্গে অধ্যয়ন করা হলে এই আকর্ষণীয় কিন্তু অত্যন্ত কল্পনাপ্রসূত তত্ত্ব টেকেনি। ব্যাবিলনীয় জ্যোতির্বিজ্ঞান সম্পর্কে আমাদের বর্তমান প্রকৃত জ্ঞান মূলত তিনজন জেসুইট পণ্ডিতের কাজের ওপর দাঁড়িয়ে আছে: অ্যাসিরিওলজিস্ট জে. এন. স্ট্রাসমায়ার এবং জ্যোতির্বিদ জে. এপিং ও এফ. এক্স. কুগলার।

প্রথম ব্যাবিলনীয় রাজ্যের যুগ ছিল অর্থনৈতিক ও রাজনৈতিক শক্তি এবং সাংস্কৃতিক জীবনের চূড়ান্ত বিকাশের সময়। বাণিজ্য ও শিল্প সমৃদ্ধ হয়েছিল। অত্যন্ত গুরুত্বপূর্ণ বাণিজ্যনগরী ব্যাবিলন কেবল একটি বিশাল মেসোপটেমীয় সাম্রাজ্যের রাজধানীই ছিল না, বরং সমগ্র নিকট-এশিয়ায় আলো ছড়ানো একটি সাংস্কৃতিক কেন্দ্রও ছিল। এখানে পূর্ববর্তী শতাব্দীগুলোর সুমেরীয় সংস্কৃতির ফল পূর্ণতা পায়, আর ধর্মতত্ত্বের ব্যবস্থা এমন রূপ নেয় যা পরবর্তীকালেও বহাল থাকে। বিখ্যাত "হাম্মুরাবির আইনসংহিতা" থেকে আমরা নাগরিক জীবনের কাঠামো ভালোভাবে জানি। এটি একটি পাথরে খোদাই করা, যা উনিশ শতকে সুসায় খনন করে পাওয়া যায়। খুঁড়ে পাওয়া অসংখ্য মাটির ফলকের মধ্যে জমি কেনাবেচা, ঋণ দেওয়া, ভাড়া ও চাকরি সংক্রান্ত বহু দেওয়ানি চুক্তি পাওয়া গেছে। এগুলো মন্দিরে জমা রেখে সংরক্ষণ করা হতো, যেন মন্দিরগুলো নোটারি অফিস। এগুলোর মাধ্যমে আমরা পঞ্জিকা সম্পর্কে এবং রাজাদের নাম ও তারিখের ক্রম সম্পর্কে নিশ্চিত হতে পেরেছি, যা ইতিহাসের কাঠামো হিসেবে কাজ করে।

উরের রাজা দুংগির আমলে প্রচলিত বারো মাসের সুমেরীয় নামগুলো বিভিন্ন জায়গায় বিভিন্ন ছিল। এগুলোর মধ্যে দেখা যায়, চতুর্থ মাসের নাম গঠিত হয়েছে "বীজ" ও "হাত"-এর চিহ্ন দিয়ে, একাদশ মাসের নাম "শস্য" ও "কাটা"-র চিহ্ন দিয়ে, আর দ্বাদশ মাসের নাম "শস্য" ও "ঘর"-এর চিহ্ন দিয়ে। এভাবে ঋতুটি স্পষ্টভাবে নির্দেশিত হয়েছে। তাই আমরা বুঝতে পারি কেন সেখানে অধিমাস যোগ করার জন্য একাদশ মাসটি দ্বিগুণ করা হতো। ব্যাবিলনের উত্থানের পর সেমিটিক নামগুলো চালু হয়: নিসান্নু, আইয়ারু, সিমান্নু, দুজু, আবু, উলুলু, তিশরিতু, আরাখ-সাম্না, কিসলিমু, তেবিতু, শাবাতু, আদারু। জ্যোতির্বৈজ্ঞানিক লেখায় এগুলোকে একটি করে চিহ্ন দিয়ে নির্দেশ করা হয়, যা তাদের আগের সুমেরীয় নামের প্রথম অক্ষর।

প্রাচীন ব্যাবিলনীয় রাজ্যে অধিমাসটি ছিল বছরের শেষে একটি দ্বিতীয় আদারু। দেওয়ানি চুক্তির তারিখগুলোতে এই অধিমাস যোগের অনিয়ম দেখা যায়। এ থেকে বোঝা যায় যে এগুলো অভিজ্ঞতার ভিত্তিতে করা হতো, ফসল পাকার অবস্থা অনুযায়ী বা অন্য কোনো কারণে প্রয়োজন মনে হলে। পঞ্জিকা অনেক বেশি পিছিয়ে পড়লে এমনকি পরপর দুই বছরও ১৩ মাসের হয়েছে। কয়েকবার, স্পষ্টতই কোনো বিচ্যুতি দ্রুত সংশোধনের প্রয়োজন হলে, ষষ্ঠ মাসটি দ্বিগুণ করা হয়েছিল। এমন একটি ঘটনার দলিল আছে, যাতে লেখা: "হাম্মুরাবি এভাবে বলছেন: 'যেহেতু বছরটি ভালো নয়, তাই পরের মাসটিকে দ্বিতীয় উলুলু হিসেবে লিখতে হবে। তিশরিতুর ২৫ তারিখে ব্যাবিলনে দশমাংশ পৌঁছে দেওয়ার বদলে দ্বিতীয় উলুলুর ২৫ তারিখে পৌঁছে দেওয়ার ব্যবস্থা করো।'" রাজদরবারের জন্য খাদ্য সরবরাহ অবশ্যই এক মাস পিছিয়ে দেওয়া যেত না।

পঞ্জিকার জন্য কিছু আকাশের ঘটনা পর্যবেক্ষণ করা জরুরি ছিল। মাসের প্রথম দিনটি সঠিকভাবে নির্ধারণ করতে নতুন চাঁদকে তার প্রথম আবির্ভাবেই ধরতে হতো। এটি খুব কঠিন ছিল না। সেই অপূর্ব উজ্জ্বল জলবায়ুতে (শীতের কয়েকটি মাস আর মাঝে মাঝে বালুঝড় বাদে) বিস্তীর্ণ সমভূমির এক দিগন্ত থেকে আরেক দিগন্ত পর্যন্ত স্পষ্ট দেখা যেত। পুরোহিত-জ্যোতির্বিদেরা তাঁদের ধাপযুক্ত মিনার থেকে সান্ধ্য আকাশে চাঁদের প্রথম সরু ফালি সহজেই দেখতে পেতেন। আনুষ্ঠানিক প্রয়োজনে তাঁদের পূর্ণিমার দিকেও নজর রাখতে হতো। আর পূর্বাভাসের জন্য মাসের শেষ দিকে প্রভাতী বাঁকা চাঁদের শেষ আবির্ভাবও পর্যবেক্ষণ করতে হতো। এসব করতে গিয়ে তাঁরা কি নীরবে নিজেদের পথ আঁকতে থাকা আর প্রতি মাসে এগিয়ে যাওয়া নক্ষত্রগুলো লক্ষ করেননি, এমন কি কল্পনা করা যায়? কিংবা তাদের মাঝে মাঝে মাঝে দেখা দেওয়া উজ্জ্বল গ্রহগুলো, বিশেষ করে অতুলনীয় সন্ধ্যাতারা, কি তাঁদের মুগ্ধ করেনি? এগুলো নিয়ে লেখা খুব কম। অত পরিষ্কার আকাশ থাকা সত্ত্বেও আদি ব্যাবিলনীয়দের মন নক্ষত্রে ততটা নিবিষ্ট ছিল না, যতটা প্রায়ই ধারণা করা হয়েছে। তবে খুবই সম্ভব, এবং কিছু লেখাও এটি সমর্থন করে, যে চাঁদের এই নিয়মিত পর্যবেক্ষণ ধীরে ধীরে নক্ষত্র নিয়ে ব্যবহারিক আগ্রহ বাড়িয়ে তুলেছিল।

নতুন চাঁদ পর্যবেক্ষণে প্রথম যে ঘটনাটি ধরা পড়ে, তা হলো বছরের পরিক্রমায় তারামণ্ডলগুলোর নিয়মিত অগ্রসর হওয়া ও সরে যাওয়া। পশ্চিমের সান্ধ্য আকাশে দৃশ্যমান নক্ষত্রগুলো ঋতুর বৈশিষ্ট্য বহন করে। তাই এগুলো দিয়ে পঞ্জিকা ও অধিমাস যাচাই করা যায়। পূর্বের প্রভাতী আকাশে নক্ষত্রের প্রথম আবির্ভাবের ক্ষেত্রেও একই কথা খাটে। তাই নক্ষত্রের নিয়মিত সূর্যোদয়কালীন উদয় ও অস্তের মাধ্যমে অধিমাস যোগের পদ্ধতি ধীরে ধীরে ফসল দেখে করা অনিয়মিত অভিজ্ঞতাভিত্তিক পদ্ধতির জায়গা নিয়েছিল নিশ্চয়ই। এর একটি সুস্পষ্ট ইঙ্গিত পাওয়া যায় একটি তালিকায় (অজানা কোনো শতাব্দীর একটি মূল লেখার আসিরীয় অনুলিপি)। এতে নক্ষত্র বা তারামণ্ডলের ৩৬টি নাম আছে, প্রতি মাসের জন্য তিনটি করে, যার প্রতিটির প্রথমটি স্পষ্টতই সেই মাসের সূর্যোদয়কালীন উদয়ের সঙ্গে সম্পর্কিত। নিসান্নু মাসের জন্য পাই দিলগান (অর্থাৎ মেষ ও তিমি)। আইয়ারুর জন্য প্রথম নামটি মুলমুল (কৃত্তিকা; "মুল" মানে "নক্ষত্র")। সিমান্নুর জন্য সিবজিয়ান্নু (কালপুরুষ), এবং এভাবে বাকিগুলো। এ ছাড়া আরেকটি অনেকটা ক্ষতিগ্রস্ত লেখায় লিপিবদ্ধ আছে: "দিলগান নক্ষত্র নিসান্নু মাসে দেখা দেয়; নক্ষত্রটি অনুপস্থিত থাকলে মাসটিকে অবশ্যই… মুলমুল নক্ষত্র আইয়ারু মাসে দেখা দেয়…" নক্ষত্রের আবির্ভাব বা আবির্ভাবে দেরি স্পষ্টতই নির্দিষ্ট কিছু কাজের নির্দেশনা দেয়।

